\documentclass[11pt, letterpaper]{article}

% --- Idioma y codificación ---
\usepackage[spanish]{babel}
\usepackage[utf8]{inputenc}
\usepackage[T1]{fontenc}
\usepackage{lmodern}

% --- Matemáticas ---
\usepackage{amsmath, amssymb, amsthm}

% --- Formato ---
\usepackage[margin=2.5cm]{geometry}
\usepackage{parskip}
\usepackage{booktabs}
\usepackage{hyperref}
\hypersetup{colorlinks=true, linkcolor=blue!70!black, urlcolor=blue!70!black}

% --- Entornos de teoremas ---
\theoremstyle{definition}
\newtheorem{definition}{Definición}[section]
\theoremstyle{plain}
\newtheorem{theorem}{Teorema}[section]
\newtheorem{lemma}[theorem]{Lema}
\newtheorem{corollary}[theorem]{Corolario}
\theoremstyle{remark}
\newtheorem{remark}{Observación}[section]

\title{%
  T03 --- Demostración formal:\\[4pt]
  Quicksort aleatorizado es $O(n \log n)$ en esperanza
}
\author{Curso Linux--C--Rust --- B15 Estructuras de datos en C}
\date{}

\begin{document}
\maketitle
\tableofcontents

% ===================================================================
\section{Modelo}
% ===================================================================

\textbf{Entrada.} Un array fijo $A[0..n{-}1]$ de $n$ elementos
distintos (la entrada \textbf{no} es aleatoria).

\textbf{Algoritmo.} Quicksort con pivote aleatorio: en cada llamada
recursiva sobre un subarreglo de tamaño~$m$, se elige un pivote
uniformemente al azar entre los $m$ elementos del subarreglo.

\begin{remark}
La aleatoriedad es del algoritmo, no de la entrada. A diferencia del
análisis de caso promedio (T02), aquí la entrada es fija y arbitraria.
La esperanza se toma sobre las elecciones aleatorias del pivote.
\end{remark}

% ===================================================================
\section{Distinción clave: caso promedio vs aleatorizado}
% ===================================================================

\begin{table}[h]
\centering
\begin{tabular}{lp{4.5cm}p{4.5cm}}
\toprule
& \textbf{Caso promedio (T02)} & \textbf{Aleatorizado} \\
\midrule
Entrada
  & Permutación uniforme aleatoria
  & \textbf{Fija, arbitraria} \\
Pivote
  & Determinístico (ej: último)
  & \textbf{Aleatorio uniforme} \\
Fuente de aleatoriedad
  & Los datos
  & El algoritmo \\
Garantía
  & ``En promedio sobre entradas''
  & ``Para toda entrada, en esperanza'' \\
¿Entrada adversarial?
  & Sí ($O(n^2)$ determinístico)
  & \textbf{No} \\
\bottomrule
\end{tabular}
\end{table}

La conclusión numérica es la misma ($\approx 1.39 \cdot n\log_2 n$
comparaciones esperadas), pero la garantía es fundamentalmente más
fuerte en el caso aleatorizado.

% ===================================================================
\section{Demostración 1: la esperanza es independiente de la entrada}
% ===================================================================

\begin{theorem}\label{thm:independent}
Sea $A$ cualquier array fijo de $n$ elementos distintos. El número
esperado de comparaciones de quicksort aleatorizado sobre~$A$ es:
\[
  E[C] = 2(n+1)H_n - 4n \approx 1.386 \cdot n \log_2 n
\]
\end{theorem}

\begin{proof}
Sean $z_1 < z_2 < \ldots < z_n$ los elementos de $A$ en orden
creciente (los valores, no las posiciones). Definimos:
\[
  X_{ij} = \begin{cases}
    1 & \text{si } z_i \text{ y } z_j \text{ se comparan} \\
    0 & \text{si no}
  \end{cases}
\]

Como en T02:
\[
  E[C] = \sum_{i=1}^{n-1}\sum_{j=i+1}^{n} \Pr[X_{ij} = 1]
\]

\textbf{Afirmación.} $\Pr[X_{ij} = 1] = 2/(j - i + 1)$, y esta
probabilidad \textbf{no depende de la disposición de los elementos
en~$A$}, solo de sus valores relativos.

\emph{Prueba de la afirmación.} Consideremos
$Z_{ij} = \{z_i, z_{i+1}, \ldots, z_j\}$. Independientemente de
dónde estén estos elementos en el array:
\begin{itemize}
  \item Mientras ningún elemento de $Z_{ij}$ sea elegido como pivote,
    todos permanecen en el mismo subarreglo (cualquier pivote externo
    a $Z_{ij}$ los mantiene juntos, pues son todos mayores o todos
    menores que él).
  \item El primer elemento de $Z_{ij}$ elegido como pivote determina
    si $z_i$ y $z_j$ se comparan. Como el pivote se elige
    uniformemente dentro del subarreglo, y los elementos de $Z_{ij}$
    están en el mismo subarreglo al momento de la elección, cada uno
    tiene probabilidad $1/(j{-}i{+}1)$ de ser el primero elegido.
  \item Solo si se elige $z_i$ o $z_j$ (2 opciones de $j{-}i{+}1$)
    ocurre la comparación.
\end{itemize}

Esta probabilidad depende solo de la estructura de rangos, no de las
posiciones iniciales en~$A$. Por lo tanto:
\[
  \Pr[X_{ij} = 1] = \frac{2}{j-i+1}
\]
para \textbf{cualquier} disposición de los elementos.

El resto de la derivación es idéntico a T02:
$E[C] = 2(n+1)H_n - 4n$.
\end{proof}

% ===================================================================
\section{Demostración 2: concentración}
% ===================================================================

\begin{theorem}\label{thm:concentration}
Para quicksort aleatorizado sobre cualquier entrada de tamaño~$n$:
\[
  \Pr[C > \alpha \cdot n \ln n] \leq n^{2 - \alpha}
  \quad \text{para } \alpha > 2
\]
En particular, $\Pr[C > 4n\ln n] < 1/n^2$.
\end{theorem}

\begin{proof}[Prueba (sketch usando desigualdad de Markov sobre $2^C$)]
Definimos la variable aleatoria $Y = 2^{C/(cn)}$ para una constante
$c$ adecuada. Se puede demostrar que $E[Y] = O(n)$. Por la
desigualdad de Markov:
\[
  \Pr[C > \alpha \cdot n \ln n]
  = \Pr\!\left[Y > 2^{\alpha \ln n / c}\right]
  \leq \frac{E[Y]}{2^{\alpha \ln n / c}}
\]

Con la elección adecuada de $c$, esto da una cota polinomialmente
pequeña en~$n$.
\end{proof}

\begin{remark}[Interpretación práctica]
Para $n = 10^6$:
\begin{itemize}
  \item $E[C] \approx 1.39 \cdot 10^6 \cdot 20 \approx 2.8 \times 10^7$.
  \item $\Pr[C > 4 \cdot 10^6 \cdot \ln(10^6)] < 10^{-12}$.
\end{itemize}
La probabilidad de que quicksort aleatorizado sea significativamente
más lento que su esperanza es infinitesimal.
\end{remark}

% ===================================================================
\section{Demostración 3: no existe entrada adversarial}
% ===================================================================

\begin{theorem}\label{thm:no-adversary}
Para todo array $A$ de $n$ elementos:
\[
  E_{\mathrm{rand}}[C(A)] = 2(n+1)H_n - 4n
\]
La esperanza es \textbf{la misma} para toda entrada~$A$.
\end{theorem}

\begin{proof}
Ya demostrada en el Teorema~\ref{thm:independent}: la probabilidad
$\Pr[X_{ij} = 1] = 2/(j{-}i{+}1)$ no depende de~$A$, solo de~$n$.
Por lo tanto, la esperanza $E[C]$ tampoco depende de~$A$.
\end{proof}

\begin{corollary}
Un adversario que conoce el algoritmo pero no las elecciones aleatorias
no puede construir una entrada que fuerce más de $2(n{+}1)H_n - 4n$
comparaciones en esperanza.
\end{corollary}

\begin{table}[h]
\centering
\caption{Contraste con quicksort determinístico.}
\begin{tabular}{lcc}
\toprule
\textbf{Propiedad}
  & \textbf{Determinístico}
  & \textbf{Aleatorizado} \\
\midrule
$\exists$ entrada con $T(n) = \Theta(n^2)$?
  & \textbf{Sí} (ej: ordenada)
  & No (en esperanza) \\
$E[C]$ depende de la entrada?
  & Sí
  & \textbf{No} \\
Garantía para peor entrada
  & $\Theta(n^2)$
  & $O(n\log n)$ esperado \\
$\Pr[C = \Omega(n^2)]$
  & 1 para entrada adversarial
  & $< 1/n!$ \\
\bottomrule
\end{tabular}
\end{table}

% ===================================================================
\section{Demostración 4: profundidad de recursión esperada
  \texorpdfstring{$O(\log n)$}{O(log n)}}
% ===================================================================

\begin{theorem}\label{thm:depth}
La profundidad de recursión esperada de quicksort aleatorizado es
$O(\log n)$.
\end{theorem}

\begin{proof}
En cada llamada recursiva, el pivote cae en la posición~$k$
uniformemente entre 1 y~$m$ (tamaño del subarreglo). La probabilidad
de que el pivote caiga en el ``centro'' (entre la posición $m/4$ y
$3m/4$) es $1/2$.

Cuando el pivote cae en el centro, ambas particiones tienen tamaño
$\leq 3m/4$. Así, el tamaño del subproblema se reduce por un factor
$\leq 3/4$ con probabilidad $1/2$ en cada nivel.

El número de reducciones exitosas necesarias para llegar a tamaño~1
es:
\[
  (3/4)^d \cdot n \leq 1
  \quad\Longrightarrow\quad
  d \geq \log_{4/3} n
  = \frac{\log_2 n}{\log_2(4/3)}
  \approx 2.41 \log_2 n
\]

Cada reducción exitosa ocurre con probabilidad $1/2$. La profundidad
total (incluyendo intentos fallidos) tiene esperanza:
\[
  E[\text{profundidad}] \approx 2 \cdot 2.41 \log_2 n
  \approx 4.82 \log_2 n = O(\log n)
\]

Más precisamente, la profundidad excede $c\log n$ con probabilidad
$\leq 1/n$ para $c$ suficientemente grande.
\end{proof}

% ===================================================================
\section{Resumen de resultados}
% ===================================================================

\begin{table}[h]
\centering
\begin{tabular}{llp{5.5cm}}
\toprule
\textbf{Resultado} & \textbf{Técnica} & \textbf{Clave} \\
\midrule
$E[C] = 2(n{+}1)H_n - 4n$ para toda entrada
  & Variables indicadoras
  & $\Pr[X_{ij}{=}1]$ no depende de la disposición \\
$\Pr[C > 4n\ln n] < 1/n^2$
  & Desigualdad de Markov
  & Colas exponencialmente pequeñas \\
No existe entrada adversarial
  & $E[C]$ independiente de $A$
  & La aleatoriedad es del algoritmo \\
Profundidad $O(\log n)$
  & Pivote ``centrado'' con prob $1/2$
  & Reducción geométrica $3/4$ por éxito \\
\bottomrule
\end{tabular}
\end{table}

\end{document}
